% !TEX root =./ECE343_Practice_main.tex
%
% L34 practice set (Maxwell's Equations): the voltage induced in a coil by a
% sinusoidal field, checking a pair of fields against Maxwell's equations,
% conservation of charge, and sizing an ideal transformer.

\fullwidth{\textbf{Learning Objective 8.3: } I know the forms of Maxwell's
equations for time-varying fields and use them to calculate electric and
magnetic fields for a variety of physical systems.
\\[4pt]\normalfont\small Show your work. Write a \textbf{Documentation} line at the
end (\textbf{None} if you did not collaborate).}

\begin{questions}

%%%%%%%%%%%%%%%%%%%%%%%%%%%%%%%%%%%%%%%%%
\question \textbf{A pickup coil.} A square loop $20$ cm on a side lies in the
$xy$ plane in a uniform field
$\mathbf{B} = \hat{\mathbf{z}}\ 0.5\sin(120\pi t)$ mT.
\begin{parts}

	\begin{minipage}[t]{0.58\linewidth}
		\part Find the amplitude of the voltage induced around the loop.
		\begin{solution}[1.2in]
		The flux is $\Phi = AB = (0.2)^2(0.5\times10^{-3})\sin(120\pi t)$ Wb.
		Faraday's law gives
		\eq{
			\oint_C \mathbf{E}\mycdot d\mathbf{l} &= -\dfrac{d\Phi}{dt} \\
			&= -(0.04)(0.5\times10^{-3})(120\pi)\cos(120\pi t) \\
			&= -(7.54\ \text{mV})\cos(120\pi t)
		}
		\end{solution}
	\end{minipage}\hfill%
	\begin{minipage}[t]{0.37\linewidth}
		\raggedleft \vspace{12pt}
		amplitude = \ansbox[1.3in]{$7.54$ mV}
	\end{minipage}

	\begin{minipage}[t]{0.58\linewidth}
		\part We replace the loop with a coil of $50$ turns of the same size. Find
		the new amplitude.
		\begin{solution}[0.8in]
		Each turn links the same flux, and the turns are in series:
		\eq{ 50\times7.54\ \text{mV} = 377\ \text{mV} }
		\end{solution}
	\end{minipage}\hfill%
	\begin{minipage}[t]{0.37\linewidth}
		\raggedleft \vspace{12pt}
		amplitude = \ansbox[1.3in]{$377$ mV}
	\end{minipage}

	\part Just after $t = 0$, $B$ is positive and increasing. Seen from $+z$,
	which way does the induced current flow in the single loop? Explain with
	Lenz's law.
	\begin{solution}[1in]
	Clockwise. Just after $t = 0$ the cosine is positive, so the circulation
	counterclockwise is negative and the current flows clockwise. A clockwise
	current makes a field along $-\hat{\mathbf{z}}$ inside the loop, which
	opposes the growing flux along $+\hat{\mathbf{z}}$.
	\end{solution}

\end{parts}

\newpage
%%%%%%%%%%%%%%%%%%%%%%%%%%%%%%%%%%%%%%%%%
\question \textbf{Checking a pair of fields.} A lossless, source-free region has
$\epsilon_r = 4$ and $\mu_r = 1$. A proposed electric field is
$\mathbf{E} = \hat{\mathbf{y}}\ 10\cos(2\pi\times10^{9}\ t - \beta z)$ V/m.
\begin{parts}

	\part Use Faraday's law to find the magnetic field $\mathbf{H}$ that must
	accompany $\mathbf{E}$, in terms of $\beta$.
	\begin{solution}[1.6in]
	$E_y$ depends only on $z$, so only the $\hat{\mathbf{x}}$ component of the
	curl survives:
	\eq{
		\nabla\times\mathbf{E} &= -\hat{\mathbf{x}}\ \dfrac{\partial E_y}{\partial z}
		= -\hat{\mathbf{x}}\ 10\beta\sin(\omega t - \beta z) = -\mu\dfrac{\partial\mathbf{H}}{\partial t} \\
		\mathbf{H} &= -\hat{\mathbf{x}}\ \dfrac{10\beta}{\mu\omega}\cos(\omega t - \beta z)
	}
	\end{solution}

	\begin{minipage}[t]{0.58\linewidth}
		\part Use the Ampere--Maxwell law to find the value of $\beta$ for which
		the pair satisfies Maxwell's equations, and the wavelength.
		\begin{solution}[1.6in]
		\eq{
			\nabla\times\mathbf{H} &= \hat{\mathbf{y}}\ \dfrac{\partial H_x}{\partial z}
			= -\hat{\mathbf{y}}\ \dfrac{10\beta^2}{\mu\omega}\sin(\omega t - \beta z) \\
			\epsilon\dfrac{\partial\mathbf{E}}{\partial t} &= -\hat{\mathbf{y}}\ 10\epsilon\omega\sin(\omega t - \beta z)
		}
		Matching gives $\beta = \omega\sqrt{\mu\epsilon} = \omega\sqrt{\epsilon_r}/c$:
		\eq{
			\beta &= \dfrac{(2\pi\times10^{9})(2)}{3\times10^{8}} = 41.9\ \text{rad/m} \\
			\lambda &= \dfrac{2\pi}{\beta} = 15.0\ \text{cm}
		}
		\end{solution}
	\end{minipage}\hfill%
	\begin{minipage}[t]{0.37\linewidth}
		\raggedleft \vspace{12pt}
		$\beta$ = \ansbox[1.3in]{$41.9$ rad/m}

		\vspace{4pt}
		$\lambda$ = \ansbox[1.3in]{$15.0$ cm}
	\end{minipage}

	\begin{minipage}[t]{0.58\linewidth}
		\part Find the amplitude of $\mathbf{H}$ and its direction.
		\begin{solution}[1in]
		\eq{
			\eta &= \sqrt{\dfrac{\mu}{\epsilon}} = \dfrac{377}{\sqrt{4}} = 188\ \Omega \\
			H_0 &= \dfrac{10}{188} = 53.1\ \text{mA/m}, \quad \text{along } -\hat{\mathbf{x}}
		}
		\end{solution}
	\end{minipage}\hfill%
	\begin{minipage}[t]{0.37\linewidth}
		\raggedleft \vspace{12pt}
		$H_0$ = \ansbox[1.3in]{$53.1$ mA/m}
	\end{minipage}

	\part Show that the pair also satisfies both of Gauss's laws.
	\begin{solution}[0.8in]
	$D_y$ depends only on $z$ and $t$, so $\nabla\mycdot\mathbf{D} = \partial D_y/\partial y = 0 = \rho_v$.
	$B_x$ depends only on $z$ and $t$, so $\nabla\mycdot\mathbf{B} = \partial B_x/\partial x = 0$.
	\end{solution}

\end{parts}

\newpage
%%%%%%%%%%%%%%%%%%%%%%%%%%%%%%%%%%%%%%%%%
\question \textbf{Conservation of charge.} In a region, the current density is
$\mathbf{J} = \hat{\mathbf{r}}\ kr$ A/m$^2$ with $k = 2$ A/m$^3$, in spherical
coordinates.
\begin{parts}

	\begin{minipage}[t]{0.58\linewidth}
		\part Find the rate of change of the charge density, $\partial\rho_v/\partial t$.
		\begin{solution}[1.2in]
		The continuity equation, with the spherical divergence from Lesson 17:
		\eq{
			\dfrac{\partial\rho_v}{\partial t} = -\nabla\mycdot\mathbf{J}
			= -\dfrac{1}{r^2}\dfrac{\partial}{\partial r}\left(r^2\ kr\right) = -3k = -6\ \text{C/(m}^3\text{s)}
		}
		\end{solution}
	\end{minipage}\hfill%
	\begin{minipage}[t]{0.37\linewidth}
		\raggedleft \vspace{12pt}
		$\partial\rho_v/\partial t$ = \ansbox[1.3in]{$-6$ C/(m$^3$s)}
	\end{minipage}

	\begin{minipage}[t]{0.58\linewidth}
		\part Find the current out of a sphere of radius $10$ cm, and show that it
		equals the rate at which the charge inside decreases.
		\begin{solution}[1.4in]
		\eq{
			\oint \mathbf{J}\mycdot d\mathbf{s} &= (kr)(4\pi r^2) = 4\pi k r^3 \\
			&= 4\pi(2)(0.1)^3 = 25.1\ \text{mA} \\
			-\dfrac{dQ}{dt} &= \dfrac{4}{3}\pi r^3 (3k) = 4\pi k r^3 = 25.1\ \text{mA}
		}
		\end{solution}
	\end{minipage}\hfill%
	\begin{minipage}[t]{0.37\linewidth}
		\raggedleft \vspace{12pt}
		current = \ansbox[1.3in]{$25.1$ mA}
	\end{minipage}

	\part Explain why Ampere's law without the displacement current could not
	describe this region.
	\begin{solution}[1in]
	The divergence of a curl is zero, so $\nabla\times\mathbf{H} = \mathbf{J}$
	would force $\nabla\mycdot\mathbf{J} = 0$. Here $\nabla\mycdot\mathbf{J} = 6$
	A/m$^3$, because charge is leaving the region. The displacement current
	$\partial\mathbf{D}/\partial t$ supplies the term whose divergence is
	$\partial\rho_v/\partial t$ and restores the balance.
	\end{solution}

\end{parts}

\newpage
%%%%%%%%%%%%%%%%%%%%%%%%%%%%%%%%%%%%%%%%%
\question \textbf{An ideal transformer.} A transformer runs from a $240$ V rms,
$50$ Hz line with $N_1 = 960$ turns on the primary and supplies $24$ V rms to an
$8\ \Omega$ load. The core has a cross-sectional area of $10\ \text{cm}^2$.
\begin{parts}

	\begin{minipage}[t]{0.58\linewidth}
		\part Find the number of secondary turns and both currents.
		\begin{solution}[1.2in]
		\eq{
			N_2 &= N_1\dfrac{V_2}{V_1} = 960\times\dfrac{24}{240} = 96 \\
			I_2 &= \dfrac{24}{8} = 3\ \text{A}, \qquad I_1 = I_2\dfrac{N_2}{N_1} = 0.3\ \text{A}
		}
		\end{solution}
	\end{minipage}\hfill%
	\begin{minipage}[t]{0.37\linewidth}
		\raggedleft \vspace{12pt}
		$N_2$ = \ansbox[1.3in]{$96$}

		\vspace{4pt}
		$I_1$ = \ansbox[1.3in]{$0.3$ A}
	\end{minipage}

	\begin{minipage}[t]{0.58\linewidth}
		\part Find the peak flux in the core and the peak flux density.
		\begin{solution}[1.3in]
		With $v_1 = N_1\ d\Phi/dt$ and $V_\text{1,pk} = 240\sqrt{2} = 339$ V:
		\eq{
			\Phi_\text{pk} &= \dfrac{V_\text{1,pk}}{N_1\omega} = \dfrac{339}{(960)(2\pi\times50)} = 1.13\ \text{mWb} \\
			B_\text{pk} &= \dfrac{\Phi_\text{pk}}{A} = \dfrac{1.13\times10^{-3}}{10\times10^{-4}} = 1.13\ \text{T}
		}
		\end{solution}
	\end{minipage}\hfill%
	\begin{minipage}[t]{0.37\linewidth}
		\raggedleft \vspace{12pt}
		$B_\text{pk}$ = \ansbox[1.3in]{$1.13$ T}
	\end{minipage}

	\part A student proposes saving copper by using $N_1 = 240$ and $N_2 = 24$,
	the same ratio. Explain what happens to the core.
	\begin{solution}[1in]
	The peak flux density scales as $1/N_1$, so it rises by a factor of $4$, to
	$4.50$ T. That is far above the $1.5$ to $1.7$ T at which a silicon-steel
	core saturates, so $\mu$ collapses, the magnetizing current grows sharply, and
	the ideal model no longer holds.
	\end{solution}

\end{parts}

\vspace{1.5em}
\noindent\textbf{Documentation:} \rule[-2pt]{0.6\linewidth}{0.4pt}

\end{questions}
